\section{Conclusion and Lessons}

We trained an RGB Diffusion Policy for ManiSkill Push-T for 50,000 updates and evaluated one fixed
EMA checkpoint over 300 held-out rollouts. The resulting baseline is \timean\% $\pm$ \tistd{}
percentage points across environment seeds, with \tipooled{} pooled successes. The result is
credible because the RGB replay is temporally aligned, controller-native, parallelism-matched, and
validated before training; the checkpoint and evaluation lineage is preserved.

The strongest completed evidence is T-I. T-II evaluation code and a falsifiable pose-conditioned
protocol were prepared, but infrastructure and time prevented the required rollouts; T-III and
T-IV were not executed. A substantial fraction of the work involved finding and correcting the
RGB replay alignment defect, matching simulator parallelism, establishing restart-safe GPU
training, and preserving a fixed-checkpoint evaluation. We prioritized this auditable baseline
over presenting unvalidated failure clusters or rushed RL curves.

The main lesson is methodological: tensor shapes, decreasing loss, and plausible videos are
insufficient without temporal provenance and closed-loop metrics. Likewise, failure labels require
pose-conditioned replication, and LLM-generated rewards require executable validation against the
simulator state. These gates define a concrete continuation path while keeping the claims in this
submission commensurate with the evidence actually obtained.
